\documentclass{situationist-preprint}

\title{Against Hedonism}
\msid{ETH-05-a}
\msfamily{ETH-05}
\msdiscipline{Ethics}
\mstype{Principal manuscript}
\mskeywords{hedonism; well-being; experience machine; Nozick; contact with reality; mental statism; status quo bias}

\begin{document}
\maketitle

\begin{abstract}
Hedonism holds that pleasure is the only intrinsic good, and in its theory of well-being that how well a life goes for the one who lives it depends only on the pleasures and pains it contains. We argue that hedonism is false, and that Nozick's experience machine shows why: well-being includes contact with reality. We restate the experience machine as a comparison between two lives that are alike in every experience and differ only in whether the experiences are veridical and the subject's aims are achieved. Idealized reflection ranks the life in contact with reality higher for the one who lives it. This is incompatible with any mental-state theory of well-being, of which hedonism is the leading example. We then answer the debunking objections that have dominated recent discussion: that the intuition reflects status quo bias, that it reflects distrust of the machine, that it reflects concern for others rather than for oneself, and that it is moral rather than prudential. Each objection is avoided by the comparative form of the case, which involves no choice, no plugging in and no one else. What remains is a judgment about two lives, and it tells against hedonism.
\end{abstract}

\section{The problem}

Hedonism about well-being is the view that a life goes well for the one who lives it to the extent that it contains pleasure and lacks pain. It has a long history, from Epicurus to \citet{Bentham1789} and \citet{Sidgwick1907}, who concluded after long examination that nothing is ultimately desirable except desirable consciousness. Hedonism about value in general adds that pleasure is the only intrinsic good. The two theses stand or fall together for present purposes: if something other than pleasure is intrinsically good for a person, then pleasure is not the only intrinsic good.

The most influential argument against hedonism is \citeauthor{Nozick1974}'s experience machine \citep[42--45]{Nozick1974}. Neuroscientists could stimulate your brain so that you would think and feel that you were writing a great novel, making a friend or reading an interesting book, while you floated in a tank. Should you plug in for life, with your experiences preprogrammed to be as pleasant as possible? Nozick held that most of us would not, and drew the conclusion that something matters to us besides how our lives feel from the inside: we want to do certain things and not just to have the experience of doing them, to be a certain sort of person, and to be in contact with a deeper reality.

The argument has been widely accepted, and for a time it was regarded as decisive. Recent work has questioned it. \citet{deBrigard2010} found that when subjects are told that they are already in a machine and asked whether they would unplug, many prefer to stay, and argued that the original reluctance to plug in reflects status quo bias rather than a judgment about well-being. \citet{Silverstein2000} and \citet{Crisp2006} argued that the intuition is explained by the instrumental value, in ordinary life, of contact with reality, and so does not show that contact is intrinsically good. \citet{Weijers2014} argued that the case is contaminated by irrelevant factors, such as fear of malfunction, and that the original scenario should be abandoned in favour of cleaner variants.

We argue that the experience machine, correctly stated, survives these objections. The correct statement is not a choice between plugging in and staying out, but a comparison between two lives. ETH-01-a uses the case in this form as an instance of its method of isolation: the two lives are alike in their experiences and differ in their contact with reality, and idealized reflective endorsement ranks them differently. We develop that use here.

\section{Framework}

\begin{definition}[Mental-state theory]\label{d:mst}
A theory of well-being is a \emph{mental-state theory} if, according to it, how well a life goes for the one who lives it supervenes on the subject's experiences, that is, on how the life is from the inside.
\end{definition}

Hedonism is a mental-state theory: pleasures and pains are experiences, or features of experiences, and on any version of hedonism, sensory or attitudinal, the hedonic value of a life is fixed by them \citep{Feldman2004}. Some non-hedonistic theories, on which well-being consists in experiences of other kinds, are mental-state theories as well.

\begin{definition}[Contact with reality]\label{d:contact}
A life is \emph{in contact with reality} to the extent that the subject's experiences as of the world are veridical, the persons she takes herself to be related to exist and are related to her as she takes them to be, and the aims she takes herself to be achieving are achieved.
\end{definition}

Contact with reality in this sense includes Nozick's three items: doing rather than seeming to do, being a person of a certain kind rather than seeming to be one, and being related to a real world. It differs from knowledge, which is a further good (ETH-01): a life can be in contact with reality without its subject knowing much about it.

\begin{definition}[Experiential twins]\label{d:twins}
Two lives are \emph{experiential twins} if they are alike in every experience, moment by moment, and differ only in their relations to the world outside the subject's experience.
\end{definition}

\begin{principle}[Prudential comparison]\label{p:comp}
If idealized reflection, attending only to how well each life goes for the one who lives it, ranks one of two lives above the other, then the first life contains more well-being.
\end{principle}

Principle~\ref{p:comp} is the prudential counterpart of the criterion of ETH-01: the good of outcomes is what idealized reflective endorsement converges on, and the good of a life for the one who lives it is what such reflection converges on when it considers the life from the point of view of its subject's interests. Idealized reflection, as there defined, is fully informed and vivid and reasons without error. The principle does not require that ordinary intuitions about the case be reliable; it requires that, where they are distorted, the distortion be identifiable and removable by idealization.

\section{The argument}

\begin{assumption}[The comparative verdict]\label{a:verdict}
There are experiential twins $L_1$ and $L_2$ such that $L_1$ is in contact with reality and $L_2$ is not, and idealized reflection, attending only to how well each goes for its subject, ranks $L_1$ above $L_2$.
\end{assumption}

The pair we have in mind is this. In $L_1$, a woman writes novels that are read and valued, has friends who love her, and raises children who flourish. In $L_2$, a woman has exactly the same experiences, but they are produced by a machine: there are no novels, no friends and no children, only the experiences as of them. Nothing in $L_2$ is ever felt to be amiss; she dies content, as the woman in $L_1$ does. The question is not which life she would choose, nor which life is better for others, but which life is better for her. Our verdict, and the verdict of almost everyone to whom the question is put in this form, is that $L_1$ is.

The verdict is confirmed by cases that involve no machine. \citet{Nagel1970} described a man betrayed by his friends, ridiculed behind his back and despised by those who treat him politely, and observed that it is a misfortune for him though he never discovers it. \citet{Kagan1994} described a businessman whose life's work is, unknown to him, a fraud, whose family and colleagues deceive him at every turn, and who dies happy; his life goes worse for him than the life of an experiential twin whose success and relationships are real. In each case, the deceived life is an experiential twin of an undeceived one, and the deceived life is judged worse for its subject.

\begin{theorem}\label{t:mst}
Every mental-state theory of well-being is false. In particular, hedonism about well-being is false.
\end{theorem}

\begin{proof}
By Assumption~\ref{a:verdict} and Principle~\ref{p:comp}, $L_1$ contains more well-being than $L_2$. By Definition~\ref{d:twins}, $L_1$ and $L_2$ are alike in every experience. A mental-state theory makes well-being supervene on experience (Definition~\ref{d:mst}), so it assigns the two lives the same well-being. Every mental-state theory therefore assigns a false ranking to $L_1$ and $L_2$. Hedonism is a mental-state theory.
\end{proof}

\begin{corollary}\label{c:contact}
Well-being includes contact with reality: a life's being in contact with reality is good, in itself, for the one who lives it.
\end{corollary}

\begin{proof}
$L_1$ and $L_2$ differ only in contact with reality (Definitions~\ref{d:contact} and~\ref{d:twins}), and $L_1$ contains more well-being. The difference in well-being cannot derive from any further benefit that contact brings, since, by Definition~\ref{d:twins}, every experiential benefit is shared, and the comparison is of the lives considered as wholes, so there are no later consequences to which the difference could be instrumental. Hence contact with reality contributes to well-being non-instrumentally.
\end{proof}

\begin{corollary}\label{c:only}
Pleasure is not the only intrinsic good.
\end{corollary}

\begin{proof}
By Corollary~\ref{c:contact}, contact with reality is intrinsically good for the one who has it, and it is not pleasure, since the lives $L_1$ and $L_2$ contain the same pleasures and differ in contact.
\end{proof}

Corollary~\ref{c:contact} does not say that contact with reality is all that matters besides pleasure, nor that it outweighs pleasure. A life of real achievement and constant pain may be worse than a pleasant life in the machine. The claim is that there is a further dimension of well-being, on which experiential twins can differ. This is the structure of the hybrid theories, on which well-being requires both pleasure and its being taken in things that are real and good \citep{Parfit1984}.

\section{Objections and replies}

\begin{objection}[Status quo bias]\label{o:sqb}
Subjects told that they are already in an experience machine tend to prefer staying in it. The reluctance to plug in therefore reflects a bias in favour of one's current situation, not a judgment about well-being \citep{deBrigard2010}.
\end{objection}

\begin{reply}
The comparative verdict of Assumption~\ref{a:verdict} involves no choice and no change of situation. Neither life is the status quo; both are considered from outside, as completed wholes. A bias in favour of the status quo cannot explain a ranking of two lives neither of which is one's own, and the ranking persists in this form. The reversed cases show that the decision to plug in or unplug is affected by bias, which is plausible, since such decisions involve the loss of everything one has. They do not show that the comparative judgment is. The companion manuscript ETH-05-b examines the empirical studies in detail.
\end{reply}

\begin{objection}[Instrumental value]\label{o:instr}
In ordinary life, contact with reality is the means to future pleasure: true beliefs guide action, real friends help us. Our attachment to contact is a habit of valuation formed by this instrumental role and carried over into cases where it has no point \citep{Silverstein2000,Crisp2006}.
\end{objection}

\begin{reply}
The genealogy may be true; it does not show that the judgment is mistaken. Many things first valued as means come to be valued for their own sake, and some of them, knowledge and friendship among them, are rightly so valued. Whether contact with reality is one of them is the question, and idealized reflection, which is fully informed about the instrumental role and its absence in the case, still ranks $L_1$ above $L_2$. If the hedonist holds that the ranking is an error, she must say what idealization would remove it. Full information and vivid imagination do not; they make the emptiness of $L_2$ more apparent.
\end{reply}

\begin{objection}[Other-regarding concern]\label{o:others}
What is wrong with $L_2$ is that the woman's novels have no readers, her friends do not exist, and her children are not raised. These are losses to the world, not to her. The ranking reflects moral or impersonal value, not well-being.
\end{objection}

\begin{reply}
The objection would succeed against a comparison in which $L_1$ contains goods for others that $L_2$ lacks, but the verdict is about what is good for her. The cases of Nagel and Kagan show the point most clearly: the betrayed man and the deceived businessman do others no less good than their undeceived twins, and what is worse in their lives is that their relationships and achievements are not what they take them to be. That is a loss to them. We judge that they have been wronged, but also that their lives have gone worse for them, and the second judgment does not depend on the first.
\end{reply}

\begin{objection}[Imaginative contamination]\label{o:contam}
In imagining the machine we cannot help imagining malfunctions, the operators' power over us and the horror of the tank. These features, not the absence of contact, explain the verdict \citep{Weijers2014}.
\end{objection}

\begin{reply}
The pair $L_1$ and $L_2$ can be specified without a machine. The deceived businessman is not in a tank, and nothing in his life malfunctions. The deception is complete and permanent, and he dies content. The verdict that his life goes worse for him than his twin's is as firm as the verdict about the machine, and it involves none of the features to which the objection appeals.
\end{reply}

\section{What the result does not establish}

The argument refutes mental-state theories of well-being, but it does not decide between their rivals. It is compatible with desire-satisfaction theories, on which the woman of $L_2$ fares worse because her desires to write, to be loved and to raise children are frustrated, and with objective-list theories, on which achievement, relationships and knowledge are good for a person independently of her desires. Nor does it determine how much contact with reality weighs against pleasure. Finally, it concerns well-being, not the good of outcomes. That the good of outcomes is plural, with components such as fairness and beauty that are not part of anyone's well-being, is argued in ETH-01-a; the present argument shows that even within well-being, the hedonist's single component is not enough.

\section{Conclusion}

Pleasure is not the only intrinsic good. Two lives alike in every experience, one in contact with reality and the other not, differ in how well they go for the one who lives them, and the life in contact with reality is better. Hence well-being does not supervene on experience, and contact with reality is a component of well-being in its own right. The experience machine establishes this when it is stated as a comparison of lives rather than as a choice, and in that form the debunking objections lose their grip.

\begin{thebibliography}{99}\small
\bibitem[Bentham(1789)]{Bentham1789} Bentham, J. (1789). \emph{An Introduction to the Principles of Morals and Legislation}. T. Payne.
\bibitem[Crisp(2006)]{Crisp2006} Crisp, R. (2006). \emph{Reasons and the Good}. Clarendon Press.
\bibitem[De Brigard(2010)]{deBrigard2010} De Brigard, F. (2010). If you like it, does it matter if it's real? \emph{Philosophical Psychology}, 23(1), 43--57.
\bibitem[Feldman(2004)]{Feldman2004} Feldman, F. (2004). \emph{Pleasure and the Good Life: Concerning the Nature, Varieties, and Plausibility of Hedonism}. Clarendon Press.
\bibitem[Kagan(1994)]{Kagan1994} Kagan, S. (1994). Me and my life. \emph{Proceedings of the Aristotelian Society}, 94, 309--324.
\bibitem[Nagel(1970)]{Nagel1970} Nagel, T. (1970). Death. \emph{No\^us}, 4(1), 73--80.
\bibitem[Nozick(1974)]{Nozick1974} Nozick, R. (1974). \emph{Anarchy, State, and Utopia}. Basic Books.
\bibitem[Parfit(1984)]{Parfit1984} Parfit, D. (1984). \emph{Reasons and Persons}. Clarendon Press.
\bibitem[Sidgwick(1907)]{Sidgwick1907} Sidgwick, H. (1907). \emph{The Methods of Ethics} (7th ed.). Macmillan.
\bibitem[Silverstein(2000)]{Silverstein2000} Silverstein, M. (2000). In defense of happiness: A response to the experience machine. \emph{Social Theory and Practice}, 26(2), 279--300.
\bibitem[Weijers(2014)]{Weijers2014} Weijers, D. (2014). Nozick's experience machine is dead, long live the experience machine! \emph{Philosophical Psychology}, 27(4), 513--535.
\end{thebibliography}

\end{document}
